\subsection{Upper integral}

DEFINITION 11. --- For every function \(f \in \mathcal{F}_{+}(T)\) , one defines the upper integral of f (with respect to the measure \(\mu\) ) to be the finite or infinite positive number

\[
\mu^ {*} (f) = \inf _ {g} \mu^ {\bullet} (g),\tag{9}
\]

where \(g\) runs over the set of lower semi-continuous functions that are \(\geqslant f\) . The notations \(\int^{*} f(t) d\mu(t)\) and \(\int^{*} f d\mu\) are also used. When \(T\) is locally compact, this definition coincides with the usual definition (Ch. V, §1, No.~1, Prop. 4). Clearly \(\mu^{\bullet}(f) \leqslant \mu^{*}(f)\) , with equality when \(f\) is lower semi-continuous. If \(A\) is a subset of \(T\) , one writes \(\mu^{*}(A)\) instead of \(\mu^{*}(\varphi_{A})\) , and this number is called the outer measure of \(A\) . The measurable sets with finite outer measure are called integrable sets, as in the case of locally compact spaces.

A function \(\mathbf{f}\) with values in a Banach space or in \(\overline{\mathbf{R}}\) such that \(\mu^{*}(|\mathbf{f}|)=0\) is said to be negligible; a set \(A\subset T\) is said to be negligible if \(\varphi_{A}\) is negligible, that is, if \(\mu^{*}(A)=0\) . The expression almost everywhere is introduced as in Ch. IV, §2, No.~3.

PROPOSITION 12. --- The function \(\mu^{*}\) is an encumbrance on T.

The properties \(a\) , \(b\) , \(c\) of Def. 1 of No.~1 are obvious. The proof of the property \(d\) is identical to that of Th. 3 of Ch. IV, §1, No.~3, on taking into account Props. 4 and 5 \(a\) ).

COROLLARY. --- A function f, with values in a Banach space or in \(\overline{R}\) , is negligible if and only if \(\mathbf{f}(t)=0\) almost everywhere.

One reduces immediately to the case of a positive function. The proof is then identical to that of Th. 1 of Ch. IV, §2, No.~3.

PROPOSITION 13. --- For every subset A of T, \(\mu^{*}(\mathbf{A})\) is the infimum of the outer measures of the open sets containing A.

The proof is identical to that of Prop. 19 of Ch. IV, §1, No.~4.

DEFINITION 12. --- Let f be a function defined on T, with values in a Banach space or in \(\overline{R}\) . One says that f is moderated for the measure \(\mu\) , or \(\mu\) -moderated, if f is zero on the complement of a countable union of integrable open sets. A subset A of T is said to be moderated if the function \(\varphi_{A}\) is moderated. The measure \(\mu\) is said to be moderated if the function 1 is \(\mu\) -moderated.

For example, since the encumbrance \(\mu^{\bullet}\) is locally bounded, every compact subset K of T is contained in an open set V such that \(\mu^{\bullet}(V)<+\infty\) ; a function that is zero outside a compact set is therefore moderated. A negligible function is moderated. The remarks following Def. 2 of Ch. V, §1, No.~2 can immediately be extended to the present context. In particular, the sum of a sequence of moderated positive functions is moderated.

Remarks. --- 1) On a Lindelöf space T (TG, IX, Appendix I, Def. 1), \(^{1}\) and in particular on a Souslin space (ibid., Cor. of Prop. 1), every measure is moderated. For, the open sets of finite measure form a covering of T, from which one can extract a countable covering of T.

\begin{enumerate}
\def\labelenumi{\arabic{enumi})}
\setcounter{enumi}{1}
\tightlist
\item
  Beware, however, that the existence of a sequence of Borel sets of finite measure for \(\mu\) , with union T, does not necessarily imply the existence of a sequence of open sets of finite measure with union T (in other words, does not imply that \(\mu\) is moderated). See Exer. 8.
\end{enumerate}

PROPOSITION 14. --- Let \(f \in \mathcal{F}_{+}(\mathrm{T})\) . If \(f\) is \(\mu\) -moderated, then \(\mu^{*}(f) = \mu^{\bullet}(f)\) ; if \(f\) is not \(\mu\) -moderated, then \(\mu^{*}(f) = +\infty\) .

If \(\mu^{*}(f) < +\infty\) , there exists a lower semi-continuous function \(g \geqslant f\) such that \(\mu^{\bullet}(g) < +\infty\) . For every \(n \in \mathbf{N}\) , let \(\mathbf{G}_{n}\) be the set of \(t \in \mathbf{T}\) such that \(g(t) > 1/n\) ; the set \(\mathbf{G}_{n}\) is open, one has \(\mu^{\bullet}(\mathbf{G}_{n}) \leqslant n\mu^{\bullet}(g) < +\infty\) , and \(f\) is zero outside the union of the \(\mathbf{G}_{n}\) : the function \(f\) is therefore moderated.

Next, let us show that \(\mu^{*}\) and \(\mu^{\bullet}\) have the same value for moderated functions. Since \(\mu^{*}\) and \(\mu^{\bullet}\) are encumbrances, it suffices to establish the relation \(\mu^{*}(f)=\mu^{\bullet}(f)\) when f is a positive function, bounded above by a constant M, and zero outside an open set G of finite measure, which we shall now do.

The measure \(\mu\) is the supremum, in \(\mathcal{M}(\mathrm{T})\) , of an increasing directed family \((\mu_i)_{i\in I}\) of measures with compact support: this follows at once from Prop. 9 of No.~8. Let \(g\) be a lower semi-continuous function on \(\mathrm{T}\) , between \(f\) and the lower semi-continuous function \(\mathrm{M}\varphi_{\mathrm{G}}\) . Set \(\nu_i = \mu - \mu_i\) ; then \(\mu^{\bullet} = \mu_i^{\bullet} + \nu_i^{\bullet}\) (No.~2, Remark 1), consequently

\[
\begin{array}{r l} \mu^ {\bullet} (g) - \mu^ {\bullet} (f) & = (\mu_ {i} ^ {\bullet} (g) - \mu_ {i} ^ {\bullet} (f)) + (\nu_ {i} ^ {\bullet} (g) - \nu_ {i} ^ {\bullet} (f)) \\ & \leqslant (\mu_ {i} ^ {\bullet} (g) - \mu_ {i} ^ {\bullet} (f)) + \nu_ {i} ^ {\bullet} (\mathrm{M} \varphi_ {\mathrm{G}}). \end{array}
\]

One has \(\nu_{i}^{\bullet}(\mathrm{M}\varphi_{\mathrm{G}}) = \mu^{\bullet}(\mathrm{M}\varphi_{\mathrm{G}}) - \mu_{i}^{\bullet}(\mathrm{M}\varphi_{\mathrm{G}})\) and \(\mu^{\bullet}(\mathrm{M}\varphi_{\mathrm{G}}) = \sup \mu_{i}^{\bullet}(\mathrm{M}\varphi_{\mathrm{G}})\) (No.~7, Prop. 6); the number \(\nu_{i}^{\bullet}(\mathrm{M}\varphi_{\mathrm{G}})\) may therefore be made arbitrarily small by a suitable choice of \(i\) . Thus everything comes down to showing that one can find, for any number \(c > 0\) and any index \(i \in I\) , a lower semicontinuous function \(g\) between \(f\) and \(\mathrm{M}\varphi_{\mathrm{G}}\) , such that \(\mu_{i}^{\bullet}(g) - \mu_{i}^{\bullet}(f) \leqslant c\) . Now, let \(L\) be the compact support of the measure \(\mu_{i}\) , and let \(\lambda\) be the measure \((\mu_{i})_{L}\) ; since \(\mu_{i}\) is concentrated on \(L\) , one has \(\mu_{i}^{\bullet}(h) = \mu_{i}^{\bullet}(h\varphi_{L}) = \lambda^{\bullet}(h_{L})\) for every function \(h \in \mathcal{F}_{+}(T)\) (No.~1, Lemma 1 and No.~2, Prop. 2); therefore

\[
\mu_ {i} ^ {\bullet} (g) - \mu_ {i} ^ {\bullet} (f) = \lambda^ {\bullet} (g _ {\mathrm{L}}) - \lambda^ {\bullet} (f _ {\mathrm{L}}).
\]

But L is compact; therefore \(\lambda^{\bullet} = \lambda^{*}\) , consequently there exists a lower semi-continuous function h defined on L, that is \(\geqslant f_{L}\) and is such that \(\lambda^{\bullet}(h) \leqslant \lambda^{\bullet}(f_{\mathrm{L}}) + c\) . Since the set L is closed in T, the function k equal to h on L, and to \(+\infty\) on T - L, is lower semi-continuous on T and is \(\geqslant f\) , and \(\lambda^{\bullet}(k_{\mathrm{L}}) = \lambda^{\bullet}(h) \leqslant \lambda^{\bullet}(f_{\mathrm{L}}) + c\) . It remains only to set \(g = \inf(k, M\varphi_{G})\) : g is lower semi-continuous, \(g \geqslant f\) and

\[
\mu_ {i} ^ {\bullet} (g) \leqslant \mu_ {i} ^ {\bullet} (k) = \lambda^ {\bullet} (k _ {\mathrm{L}}) \leqslant \lambda^ {\bullet} (f _ {\mathrm{L}}) + c = \mu_ {i} ^ {\bullet} (f) + c.
\]

COROLLARY 1. --- For a function to be negligible, it is necessary and sufficient that it be locally negligible and moderated.

COROLLARY 2. --- For a function f to be locally negligible, it is necessary and sufficient that every \(x \in T\) possess a neighborhood V such that \(f\varphi_{V}\) is negligible.

For, if this property is satisfied, \(f\varphi_{K}\) is negligible for every compact set K, and f is therefore locally negligible (No.~2, Prop. 2). Conversely, suppose that f is locally negligible, and let x be a point of T; x admits an open neighborhood V of finite measure. The function \(f\varphi_{V}\) is then locally negligible and moderated, hence is negligible.

COROLLARY 3. --- Let \(\mathbf{f}\) be a moderated function defined on \(\mathbf{T}\) . There exists a sequence \((\mathbf{K}_n)\) of pairwise disjoint compact sets, and a negligible set \(\mathbf{H}\) , such that \(\mathbf{f} = \mathbf{f}\varphi_{\mathrm{H}} + \sum_{n} \mathbf{f}\varphi_{\mathrm{K}_n}\) .

For, let G be a set that is a countable union of integrable open sets, such that f is zero outside G; then G is the union of a sequence \((K_{n})\) of pairwise disjoint compact sets and a locally negligible set H (No.~8, Prop. 11); but H is moderated, therefore negligible.

COROLLARY 4. --- If \(\mu\) and \(\nu\) are two measures on T such that \(\mu^{*} = \nu^{*}\) , then \(\mu = \nu\) .

For, the equality \(\mu^{*}=\nu^{*}\) implies that \(\mu^{\bullet}(f)=\nu^{\bullet}(f)\) for every positive function f that is moderated for \(\mu\) and \(\nu\) , hence for every positive function with compact support; it follows that \(\mu^{\bullet}=\nu^{\bullet}\) (No.~2, Prop. 2), then \(\mu=\nu\) (No.~2, Cor. of Prop. 2).

COROLLARY 5. --- If \(\mu\) is a moderated measure on T, there exists a sequence \((\mu_n)_{n \in \mathbf{N}}\) of measures with compact support such that \(\mu = \sum_{n \in \mathbf{N}} \mu_n\) .

By hypothesis, the constant function 1 is \(\mu\) -moderated. Let us apply Cor. 3 to the case \(f = 1\) ; thus, there exists a sequence \((\mathrm{K}_n)_{n \in \mathbf{N}}\) of pairwise disjoint compact subsets of T such that \(1 = \sum_{n \in \mathbf{N}} \varphi_{\mathrm{K}_n}\) \(\mu\) -almost everywhere.

Let \(\mu_{n}\) be the measure on T defined by the measure \(\mu_{\mathrm{K}_n}\) on \(\mathrm{K}_n\) (No.~3, Example 2). One knows (No.~6, Remark 2) that \(\mu_{n}\) has compact support, and that \(\mu_n^\bullet (f) = \mu^\bullet (f\varphi_{\mathrm{K}_n})\) for \(f\in \mathcal{F}_{+}(\mathrm{T})\) . Now, \(f\) is equal to \(\sum_{n\in \mathbf{N}}f\varphi_{\mathrm{K}_n}\) \(\mu\) -almost everywhere, whence

\[
\mu^ {\bullet} (f) = \sum_ {n \in \mathbf {N}} \mu^ {\bullet} (f \varphi_ {\mathrm{K} _ {n}}) = \sum_ {n \in \mathbf {N}} \mu_ {n} ^ {\bullet} (f).
\]

It follows that \(\mu = \sum_{n\in \mathbf{N}}\mu_n\) (No.~7, Prop. 7).

